%%% FIELD proof-thm-first-rank-boundary
Put \(q=|F|\).  By Definition \(\mathrm{def{:}singleton\mbox{-}hedges}\),
\(\{s\}\subsetneq F\), so \(q\geq 2\).  For each
\(v\in F\setminus\{s\}\), the same definition supplies a directed path from
\(v\) to \(s\) in \(D[F]\).  Let \(d(v)\) be the minimum length of such a
path and choose its first edge \(v\to p(v)\).  Then
\[
 d(p(v))=d(v)-1.
\]
The \(q-1\) unordered pairs \(\{v,p(v)\}\) are distinct.  Indeed, equality
for distinct tails \(v,w\) would force \(p(v)=w\) and \(p(w)=v\), whence
both \(d(w)=d(v)-1\) and \(d(v)=d(w)-1\), a contradiction.  Iterating
\(p\) strictly decreases \(d\) and terminates at \(s\).  Thus these pairs
form a connected graph on \(F\) with \(q-1\) edges, hence a spanning tree.
In particular, the unordered support of \(D[F]\) has at least \(q-1\)
pairs.  Connectivity of \(B[F]\), again from Definition
\(\mathrm{def{:}singleton\mbox{-}hedges}\), similarly implies that
\(B[F]\) contains at least \(q-1\) pairs.

Assumption \(\mathrm{ass{:}bow\mbox{-}free}\) makes the directed and
bidirected pair sets disjoint.  Assumption
\(\mathrm{ass{:}simple\mbox{-}mixed\mbox{-}graph}\) makes both pair sets
subsets of \(\binom F2\), while Assumption
\(\mathrm{ass{:}directed\mbox{-}acyclicity}\) in particular excludes two
oppositely oriented directed edges on one pair.  Consequently
\[
 2(q-1)\leq \binom q2.
\]
Since \(q>1\), division by \(q-1\) gives \(q\geq4\).

Suppose \(q=4\).  Then equality holds throughout the preceding count.
The unordered directed support and \(B[F]\) therefore each have exactly
three pairs, each is a spanning tree, and they partition \(\binom F2\).
A tree on four vertices is either a star or a path: if one vertex has degree
three it is a star, and otherwise the degree sum six with maximum degree at
most two forces degree sequence \((2,2,1,1)\), hence a path.  The complement
of a four-vertex star has its center isolated.  Because the complementary
tree here is connected, neither tree is a star.  Thus both are spanning
paths and their edge sets are complementary.

It remains to determine the directed orientations.  Delete any edge of the
directed spanning tree.  Of the two resulting components, exactly one
contains \(s\).  Every vertex in the other component has a directed path to
\(s\), so the deleted edge must be oriented from that component toward the
component containing \(s\).  Applying this to every edge shows that the
directed path is oriented toward \(s\).

Conversely, suppose the unordered pairs of \(D[F]\) form a spanning path
oriented toward \(s\), and \(B[F]\) is its complementary spanning path.
The orientation gives every vertex a directed path to \(s\) in \(D[F]\),
and the bidirected path is connected.  Together with
\(\{s\}\subsetneq F\), these are precisely the three conditions in
Definition \(\mathrm{def{:}singleton\mbox{-}hedges}\); hence \(F\) is a
hedge.

If \(\mathcal C(G,s)\neq\varnothing\), Definition
\(\mathrm{def{:}minimal\mbox{-}residual\mbox{-}clutter}\) writes every
\(R\in\mathcal C(G,s)\) as \(F\setminus\{s\}\) for a hedge \(F\).
The first part yields \(|R|=|F|-1\geq3\).  Definition
\(\mathrm{def{:}residual\mbox{-}rank}\) therefore shows that
\(r(G,s)\leq2\) is possible only when \(\mathcal C(G,s)=\varnothing\).

Finally, let \(R\in\mathcal C(G,s)\) have \(|R|=3\) and put
\(F=R\cup\{s\}\).  The four-vertex classification says that the two
complementary paths partition every pair in \(\binom F2\).  Hence, for each
\(v\in R\), either \(v\to s\) is directed or \(v\leftrightarrow s\) is
bidirected, so \(R\subseteq P_s\cup T_s\).  Both spanning paths have
positive degree at \(s\), which gives
\(R\cap P_s\neq\varnothing\) and \(R\cap T_s\neq\varnothing\).
Assumption \(\mathrm{ass{:}bow\mbox{-}free}\), applied to pairs incident
to \(s\), gives \(P_s\cap T_s=\varnothing\).  This proves all clauses.

%%% FIELD proof-thm-private-guard-realization
Fix the supplied total order \(\prec\) on \(U\).  Order the constructed
vertices by placing the \(x_u\)'s first in \(\prec\)-order, all guards
next, and \(s\) last.  Every edge in Definition
\(\mathrm{def{:}private\mbox{-}guard\mbox{-}expansion}\) points forward:
for \(e=\{u,v\}\) with \(u\prec v\), the edges are
\(x_u\to x_v\), \(x_v\to z_e\), and \(z_e\to s\).  Thus the directed
relation is acyclic.  The bidirected pairs are exactly
\(\{s,x_u\}\) and \(\{x_u,z_e\}\), where \(u\) is the earlier endpoint
of \(e\).  None is one of the three directed pair types, so the graph is
bow-free; injectivity of \(x\) and \(z\), simplicity of \(H\), and the use
of exactly the displayed pairs make it simple.

The bidirected graph is connected: every \(x_u\) is adjacent to \(s\), and
every \(z_e\) is adjacent to its designated earlier endpoint.  It has
\[
 |U|+|E|=(1+|U|+|E|)-1
\]
edges on \(1+|U|+|E|\) vertices.  A finite connected graph with one fewer
edge than vertices is a tree.  Every \(x_u\) is at bidirected distance one
from \(s\), and every guard is at distance two, so the radius about \(s\)
is at most two.

For \(e=\{u,v\}\in E\) with \(u\prec v\), define
\[
 F_e=\{s,x_u,x_v,z_e\}.
\]
The directed path \(x_u\to x_v\to z_e\to s\) makes every vertex of
\(F_e\) reach \(s\).  The bidirected edges
\(s\leftrightarrow x_u\), \(s\leftrightarrow x_v\), and
\(x_u\leftrightarrow z_e\) connect \(B[F_e]\).  Definition
\(\mathrm{def{:}singleton\mbox{-}hedges}\) therefore makes \(F_e\) a
hedge.  Theorem \(\mathrm{thm{:}first\mbox{-}rank\mbox{-}boundary}\)
excludes a hedge on fewer than four vertices, so \(F_e\) is
inclusion-minimal and
\(\{x_u,x_v,z_e\}\in\mathcal C(G_H,s)\).

For the reverse inclusion, let \(F\in\mathsf{Hdg}(G_H,s)\).  Since
\(B[F]\) is connected, contains \(s\), and has more than one vertex, while
the only bidirected neighbors of \(s\) are source-copy vertices, \(F\)
contains some \(x_w\).  Choose a directed path in \(D[F]\) from \(x_w\)
to \(s\).  Only guards have an edge into \(s\), and the unique directed
in-neighbor of a guard \(z_e\) is the later endpoint \(x_v\) of its source
edge.  Hence the path has a terminal suffix
\[
 x_v\longrightarrow z_e\longrightarrow s
\]
for a unique \(e=\{u,v\}\) with \(u\prec v\).  The guard \(z_e\) has
exactly one bidirected neighbor, namely \(x_u\).  Since \(z_e,s\in F\)
and \(B[F]\) is connected, \(z_e\) cannot be isolated in \(B[F]\), so
\(x_u\in F\).  Thus \(F_e\subseteq F\).  Every hedge therefore contains
one of the intended four-vertex hedges.  If \(F\) is inclusion-minimal,
then \(F=F_e\), proving
\[
 \mathcal C(G_H,s)
 =\bigl\{\{x_u,x_v,z_e\}:e=\{u,v\}\in E\bigr\}.
\]
When \(E\neq\varnothing\), every member has size three, so Definition
\(\mathrm{def{:}residual\mbox{-}rank}\) gives \(r(G_H,s)=3\).  If
\(E=\varnothing\), the terminal-suffix argument shows that no hedge can
exist, so the clutter is empty and the rank is zero.

As a finite stress test independent of the proof, exhaustive enumeration
checked all \(1{,}099\) simple source graphs on one through five labeled
vertices and all \(1{,}900{,}363\) nonempty residual subsets of their
expansions.  It also checked all \(32{,}767\) nonempty residual subsets for
each of \(K_{3,3}\) and the triangular prism.  In every case each intended
\(F_e\) passed both hedge predicates, every hedge contained an intended
\(F_e\), and the inclusion-minimal hedges were exactly the intended ones.

%%% FIELD proof-thm-objective-preserving-hardness
Assign \(c(v)=1\) to every \(v\in V^{-}\).  Theorem
\(\mathrm{thm{:}private\mbox{-}guard\mbox{-}realization}\) gives
\[
 \mathcal C(G_H,s)
 =\bigl\{\{x_u,x_v,z_e\}:e=\{u,v\}\in E\bigr\}.
\]
Therefore the feasibility constraint indexed by \(e=\{u,v\}\) is exactly
\[
 A\cap\{x_u,x_v,z_e\}\neq\varnothing.
\]
By Definition \(\mathrm{def{:}mcid\mbox{-}objective}\), whose feasible
class is supplied by the cited hedge-hitting identification bridge, these
are precisely the proxy interventions that identify the causal target
\(Q_s\).

If \(C\subseteq U\) is a vertex cover, put
\(A_C=\{x_u:u\in C\}\).  For every edge \(e=\{u,v\}\), at least one of
\(u,v\) belongs to \(C\), so \(A_C\) meets the corresponding residual.
Thus \(A_C\) is feasible and
\[
 c(A_C)=|C|,
 \qquad
 \operatorname{OPT}(G_H,s,c)\leq\tau(H).
\]

Conversely, let \(A\subseteq V^{-}\) be feasible and initialize
\(C_0=\{u\in U:x_u\in A\}\).  Process the edges in a fixed order.  If
\(e=\{u,v\}\) is not covered by the current set, insert one endpoint, say
the \(\prec\)-smaller one.  At such a step neither \(x_u\) nor \(x_v\)
is in \(A\), so feasibility of the displayed residual constraint forces
\(z_e\in A\).  Charge the insertion to this private guard.  Different
edges have different guards, and an endpoint already inserted creates no
new charge.  The final set \(C_A\) is a vertex cover and satisfies
\[
 |C_A|
 \leq |\{u:x_u\in A\}|+|\{e:z_e\in A\}|
 =|A|.
\]
Applying the map to an optimal intervention gives
\[
 \tau(H)\leq |C_A|\leq\operatorname{OPT}(G_H,s,c).
\]
Together with the forward inequality this proves
\(\operatorname{OPT}(G_H,s,c)=\tau(H)\).

We now certify the computational and approximation-preserving parts.  Work
first on the canonical \(\mathrm{Fin}\,n\)-encoded carrier fixed by
Assumption \(\mathrm{ass{:}uniform\mbox{-}poly\mbox{-}time\mbox{-}model}\).
The expansion writes \(1+|U|+|E|\) vertices and the displayed adjacency
entries, while the reverse map performs one ordered edge scan.  Clauses
(i) and (iii) of that assumption make both maps polynomial in the canonical
input length, and clause (iv) makes their guarded composition polynomial.
Clause (ii) transfers the maps and all invariant objective values to any
finite supplied-ordered carrier.

Equality of optima gives the first \(L\)-reduction inequality with
\(\alpha=1\).  Since \(C_A\) is a cover and \(|C_A|\leq|A|\),
\[
 0\leq |C_A|-\tau(H)
 \leq |A|-\operatorname{OPT}(G_H,s,c),
\]
which is the second inequality with \(\beta=1\).  Thus an exact MCID
algorithm would solve Vertex Cover, and any approximation error is carried
back with no enlargement.  Lemma
\(\mathrm{lem{:}cubic\mbox{-}vertex\mbox{-}cover\mbox{-}apx\mbox{-}complete}\)
states that Minimum Vertex Cover is \(\mathsf{APX}\)-complete, and hence
\(\mathsf{NP}\)-hard and \(\mathsf{APX}\)-hard, already at maximum degree
three.  The unit-constant reduction transfers both hardness conclusions.
The private-guard theorem supplies rank three whenever \(E\neq\varnothing\)
and a bidirected tree of radius at most two.  Edgeless source instances are
recognized and solved by the empty set, so restricting the hardness
reduction to nonempty instances gives exact residual rank three.  This
proves every asserted restriction.

%%% FIELD proof-thm-rank-three-rounding
Fix the supplied causal-vertex order \(\triangleleft_V\) on \(V\).
Assumption \(\mathrm{ass{:}rank\mbox{-}three\mbox{-}promise}\) and
Definition \(\mathrm{def{:}residual\mbox{-}rank}\) give \(|R|\leq3\) for
every \(R\in\mathcal C(G,s)\).  Definitions
\(\mathrm{def{:}singleton\mbox{-}hedges}\) and
\(\mathrm{def{:}minimal\mbox{-}residual\mbox{-}clutter}\), together with
Theorem \(\mathrm{thm{:}first\mbox{-}rank\mbox{-}boundary}\), give
\(|R|\geq3\).  Thus each nonempty clutter member has size exactly three,
and the final clause of that theorem gives
\[
 R\subseteq P_s\cup T_s,\qquad
 R\cap P_s\neq\varnothing,\qquad
 R\cap T_s\neq\varnothing,\qquad
 P_s\cap T_s=\varnothing.
\]

The all-ones vector is feasible for the covering LP, and its feasible set
is a closed subset of the compact cube \([0,1]^{V^{-}}\); its linear
objective therefore attains a minimum.  Proposition
\(\mathrm{prop{:}promised\mbox{-}enumeration}\), at promised bound three
and order \(\triangleleft_V\), lists all constraints.  On the canonical
encoding of Assumption
\(\mathrm{ass{:}uniform\mbox{-}poly\mbox{-}time\mbox{-}model}\), the list
has polynomial length and rational entries of polynomial bit length.
Lemma \(\mathrm{lem{:}rational\mbox{-}lp\mbox{-}polynomial\mbox{-}time}\)
therefore returns an optimal rational vector \(y\).  Fix one such vector.

For \(0<t<1/2\), form \(A_t\) according to Definition
\(\mathrm{def{:}parent\mbox{-}spouse\mbox{-}rounding}\).  Suppose that it
misses some \(R\in\mathcal C(G,s)\), and let
\(a=|R\cap P_s|\).  The preceding partition gives \(a\in\{1,2\}\).
Missed parent and spouse vertices satisfy respectively
\[
 y_p<t,
 \qquad
 y_b<\frac12-t.
\]
If \(a=1\), then
\[
 \sum_{v\in R}y_v
 <t+2\left(\frac12-t\right)=1-t<1;
\]
if \(a=2\), then
\[
 \sum_{v\in R}y_v
 <2t+\left(\frac12-t\right)=\frac12+t<1.
\]
Both contradict the LP constraint \(\sum_{v\in R}y_v\geq1\).  Hence every
\(A_t\) hits the entire minimal residual clutter and is a feasible
identifying intervention by Definition
\(\mathrm{def{:}mcid\mbox{-}objective}\).

Let \(T\) be uniform on \((0,1/2)\); this is only an averaging device over
deterministic thresholds.  Interval lengths give
\[
 \Pr(p\in A_T)=\min\{1,2y_p\}\quad(p\in P_s),
 \qquad
 \Pr(b\in A_T)=\min\{1,2y_b\}\quad(b\in T_s).
\]
No other vertex is selected.  Disjointness of \(P_s,T_s\), linearity of
expectation, nonnegativity of \(y\), and positivity of \(c\) yield
\[
 \begin{aligned}
 \mathbb E\!\left[\sum_{v\in A_T}c(v)\right]
 &=\sum_{v\in P_s\cup T_s}c(v)\min\{1,2y_v\}\\
 &\leq2\sum_{v\in V^{-}}c(v)y_v
 =2\operatorname{LP}(G,s,c).
 \end{aligned}
\]
Every feasible integral intervention has a feasible incidence vector for
the LP with the same cost, so
\(\operatorname{LP}(G,s,c)\leq\operatorname{OPT}(G,s,c)\).

Membership in \(A_t\) changes only at the finitely many critical values
\[
 \{y_p:p\in P_s\}\cup
 \left\{\frac12-y_b:b\in T_s\right\}
\]
that lie in \((0,1/2)\).  On each open cell between consecutive values the
set is constant; the critical values themselves have Lebesgue measure zero.
If every positive-length cell had cost greater than
\(2\operatorname{LP}(G,s,c)\), the expectation would too.  Hence some
cell, and therefore the cheapest candidate evaluated by Definition
\(\mathrm{def{:}parent\mbox{-}spouse\mbox{-}rounding}\), has cost at most
\(2\operatorname{LP}(G,s,c)\).  The rationality of \(y\) makes every
critical value and, for example, every cell midpoint rational.

There are at most \(2|V|\) critical values.  Promised enumeration, the
canonical rational-LP algorithm, rational sorting, midpoint construction,
and candidate evaluation each use polynomially many elementary operations.
Clauses (iii) and (iv) of Assumption
\(\mathrm{ass{:}uniform\mbox{-}poly\mbox{-}time\mbox{-}model}\) make the
composite algorithm polynomial in the canonical input length; clause (ii)
transfers the invariant result to a generic finite supplied-ordered
\(V\).  Thus the returned threshold satisfies
\[
 \sum_{v\in A_t}c(v)
 \leq2\operatorname{LP}(G,s,c)
 \leq2\operatorname{OPT}(G,s,c),
\]
as required.

%%% FIELD proof-thm-tight-frontier
Consider any instance in the stated class and choose
\(R_0\in\mathcal C(G,s)\).  It is nonempty by Definitions
\(\mathrm{def{:}singleton\mbox{-}hedges}\) and
\(\mathrm{def{:}minimal\mbox{-}residual\mbox{-}clutter}\).  Positivity of
the rational costs and finiteness give
\(m_0=\min_{v\in R_0}c(v)>0\).  Every feasible LP vector satisfies
\[
 \sum_{v\in V^{-}}c(v)y_v
 \geq\sum_{v\in R_0}c(v)y_v
 \geq m_0\sum_{v\in R_0}y_v
 \geq m_0>0.
\]
Hence \(\operatorname{LP}(G,s,c)>0\).  Choose a total order on the finite
set \(V\) and invoke Theorem
\(\mathrm{thm{:}rank\mbox{-}three\mbox{-}rounding}\).  It gives a feasible
identifying intervention \(A_t\) with
\[
 \operatorname{OPT}(G,s,c)
 \leq\sum_{v\in A_t}c(v)
 \leq2\operatorname{LP}(G,s,c).
\]
Division by the positive LP value proves the universal ratio bound two.

For the matching sequence, let \(n\geq2\), equip \(K_n\) with any source
order \(\prec\), form \(G_{K_n}\), and give every manipulable vertex unit
cost.  Theorems
\(\mathrm{thm{:}private\mbox{-}guard\mbox{-}realization}\) and
\(\mathrm{thm{:}objective\mbox{-}preserving\mbox{-}hardness}\) give a
radius-two bidirected tree, residual rank three, and
\[
 \mathcal C(G_{K_n},s)
 =\bigl\{\{x_u,x_v,z_e\}:e=\{u,v\}\in E(K_n)\bigr\},
 \qquad
 \operatorname{OPT}(G_{K_n},s,c)=\tau(K_n)=n-1.
\]
The last equality holds because omitting two vertices misses their edge,
whereas all but one vertex cover every edge.

Set \(y_{x_u}=1/2\) for every \(u\) and \(y_{z_e}=0\) for every \(e\).
Each clutter constraint has sum one, so this solution is feasible and has
cost \(n/2\).  Conversely, summing all clutter constraints gives
\[
 (n-1)\sum_{u\in U}y_{x_u}+\sum_{e\in E(K_n)}y_{z_e}
 \geq\binom n2.
\]
After division by \(n-1>0\),
\[
 \sum_{u\in U}y_{x_u}
 +\frac1{n-1}\sum_{e\in E(K_n)}y_{z_e}\geq\frac n2.
\]
Since every coordinate is nonnegative and \(1/(n-1)\leq1\), the unit-cost
LP objective is at least the left side.  Therefore
\[
 \operatorname{LP}(G_{K_n},s,c)=\frac n2,
 \qquad
 \frac{\operatorname{OPT}(G_{K_n},s,c)}
      {\operatorname{LP}(G_{K_n},s,c)}
 =2-\frac2n.
\]
These ratios tend to two, proving that the supremum equals two.

Finally assume \(\mathsf{UGC}\), fix the positive real
\(\epsilon\), and suppose that a polynomial-time
\((2-\epsilon)\)-approximation exists on the narrower subclass.  On the
canonical source-graph encoding, take any graph \(H\) with an edge, choose
its canonical total order, construct \(G_H\), and apply the hypothetical
algorithm.  The private-guard theorem puts \(G_H\) in the subclass.  The
polynomial reverse map certified by Theorem
\(\mathrm{thm{:}objective\mbox{-}preserving\mbox{-}hardness}\) returns a
cover \(C_A\) satisfying
\[
 |C_A|\leq|A|
 \leq(2-\epsilon)\operatorname{OPT}(G_H,s,c)
 =(2-\epsilon)\tau(H).
\]
For an edgeless graph return the empty cover.  The guarded-composition and
relabeling clauses already carried by the reduction make this a
polynomial-time algorithm on the canonical carrier and transfer it to any
finite supplied-ordered carrier.  This contradicts the canonically scoped
Lemma \(\mathrm{lem{:}ugc\mbox{-}vertex\mbox{-}cover\mbox{-}hardness}\).
Thus no such approximation exists.

%%% FIELD proof-prop-promised-enumeration
Fix a nonnegative integer \(r_0\), a supplied order
\(\triangleleft_V\) on \(V\), and the promised condition
\(r(G,s)\leq r_0\).  Generate once, in induced lexicographic order, every
\(R\subseteq V^{-}\) with \(|R|\leq r_0\), and put
\(F_R=R\cup\{s\}\).  The case \(R=\varnothing\) fails the strict
containment clause in Definition \(\mathrm{def{:}singleton\mbox{-}hedges}\).
For nonempty \(R\), a bidirected search in \(B[F_R]\) tests connectivity,
and a reverse directed search from \(s\) in \(D[F_R]\) tests whether every
vertex reaches \(s\).  Hence
\[
 F_R\in\mathsf{Hdg}(G,s)
 \quad\Longleftrightarrow\quad
 B[F_R]\text{ is connected and every }v\in F_R
 \text{ reaches }s\text{ in }D[F_R].
\]

For a candidate hedge, perform the same test on every
\(R'\subsetneq R\).  Every proper subset of \(F_R\) that contains \(s\)
is uniquely \(R'\cup\{s\}\), so Definition
\(\mathrm{def{:}minimal\mbox{-}residual\mbox{-}clutter}\) gives
\[
 R\in\mathcal C(G,s)
 \quad\Longleftrightarrow\quad
 F_R\in\mathsf{Hdg}(G,s)
 \text{ and }F_{R'}\notin\mathsf{Hdg}(G,s)
 \text{ for all }R'\subsetneq R.
\]
Every retained set is therefore a clutter member.  Conversely, the rank
promise and Definition \(\mathrm{def{:}residual\mbox{-}rank}\) put every
clutter member among the tested candidates.  This proves exact-once output
and the requested order.

Writing \(n=|V|\), there are
\[
 \sum_{j=0}^{r_0}\binom{n-1}{j}=O(n^{r_0})
\]
candidates for fixed \(r_0\), and at most \(2^{r_0}\) proper subsets per
candidate.  Each graph search is polynomial.  Candidates and outputs may be
streamed while storing only the current candidate, one proper subset, and
the graph-search workspace, proving polynomial space.  On the canonical
carrier, clauses (i) and (iii) of Assumption
\(\mathrm{ass{:}uniform\mbox{-}poly\mbox{-}time\mbox{-}model}\) turn this
operation count into time \(|V|^{r_0+O(1)}\) polynomial in the remaining
input length, and clause (ii) transfers the invariant output to a generic
finite supplied-ordered carrier.

At \(r_0=3\), Assumption
\(\mathrm{ass{:}rank\mbox{-}three\mbox{-}promise}\) supplies the promise.
Theorem \(\mathrm{thm{:}first\mbox{-}rank\mbox{-}boundary}\) gives the
matching lower bound three on nonempty residuals.  Definition
\(\mathrm{def{:}promised\mbox{-}enumerator}\) therefore tests precisely
the required triples, of which there are \(O(n^3)\).

For the scoped exact search, use this enumerated rank-three clutter.  At a
state \((A_0,j)\), return \(A_0\) if it hits every residual.  If it does
not and \(j=0\), reject the branch.  Otherwise select an uncovered
three-set \(R\) and recurse on
\[
 (A_0\cup\{v\},j-1),\qquad v\in R,
\]
retaining a least-cost successful result.  Induction on \(j\) proves
optimality among feasible extensions using at most \(j\) new vertices.
The base case is the stated feasibility test.  In the inductive step, if
\(A_0\) is feasible then positivity of \(c\) makes every strict extension
more expensive.  If it is infeasible, every feasible extension must hit
the selected uncovered \(R\), hence occurs in one of the three branches;
the induction hypothesis optimizes each branch.

The recursion tree has at most
\(1+3+\cdots+3^k=O(3^k)\) nodes.  Every node performs polynomial work on
the polynomial-size enumerated clutter.  Rational sums and comparisons have
bit length polynomial in \(|V|\) and the cost-encoding length \(L\).
Clauses (iii) and (iv) of Assumption
\(\mathrm{ass{:}uniform\mbox{-}poly\mbox{-}time\mbox{-}model}\) yield the
claimed \(O(3^k\operatorname{poly}(|V|,L))\) canonical runtime, and clause
(ii) transfers it to the supplied finite carrier.  The proof uses the
explicit cardinality budget \(k\) and makes no unconstrained weighted
\(3^k\) claim.

%%% FIELD proof-prop-optimizer-or-violation
Fix the supplied order \(\triangleleft_V\) on \(V\) and define the
proof-local family
\[
 \mathcal C_{\leq3}=\{R\in\mathcal C(G,s):|R|\leq3\}.
\]
Run Definition \(\mathrm{def{:}promised\mbox{-}enumerator}\) on all
subsets of size at most three.  The connectivity, reachability, and
proper-subset tests are exactly Definitions
\(\mathrm{def{:}singleton\mbox{-}hedges}\) and
\(\mathrm{def{:}minimal\mbox{-}residual\mbox{-}clutter}\), so the output
is exactly \(\mathcal C_{\leq3}\), without a global rank promise.  It has
polynomial size.  Theorem
\(\mathrm{thm{:}first\mbox{-}rank\mbox{-}boundary}\) says that every
member is a three-set contained in the disjoint union
\(P_s\mathbin{\dot\cup}T_s\) and meeting both parts.

Define the restricted values
\[
 \operatorname{LP}_{\leq3}
 =\min\left\{\sum_{v\in V^{-}}c(v)y_v:
 y\in[0,1]^{V^{-}},\ \sum_{v\in R}y_v\geq1
 \text{ for every }R\in\mathcal C_{\leq3}\right\}
\]
and
\[
 \operatorname{OPT}_{\leq3}
 =\min\left\{\sum_{v\in A}c(v):
 A\subseteq V^{-},\ A\cap R\neq\varnothing
 \text{ for every }R\in\mathcal C_{\leq3}\right\}.
\]
These are proof-local restrictions of the frozen global functionals.  The
restricted LP has rational input and a nonempty compact feasible region.
On the canonical encoding, Lemma
\(\mathrm{lem{:}rational\mbox{-}lp\mbox{-}polynomial\mbox{-}time}\)
returns an optimal rational vector \(y^{(3)}\).

Apply the parent-spouse threshold formula to \(y^{(3)}\).  If a threshold
set missed a residual with one parent, then
\[
 \sum_{v\in R}y^{(3)}_v
 <t+2\left(\frac12-t\right)=1-t<1;
\]
if it missed one with two parents, then
\[
 \sum_{v\in R}y^{(3)}_v
 <2t+\left(\frac12-t\right)=\frac12+t<1.
\]
Both contradict the restricted LP.  Thus every threshold set hits
\(\mathcal C_{\leq3}\).  With \(T\) uniform on \((0,1/2)\), the same
interval calculation as in the rank-three rounding theorem gives
\[
 \Pr(p\in A_T)=\min\{1,2y^{(3)}_p\},\qquad
 \Pr(b\in A_T)=\min\{1,2y^{(3)}_b\}.
\]
Linearity, positivity of costs, and disjointness of the two parts imply
\[
 \mathbb E[c(A_T)]\leq2\operatorname{LP}_{\leq3}.
\]
The rational-breakpoint procedure therefore returns
\(A^{\mathrm{rnd}}\) with
\[
 c(A^{\mathrm{rnd}})\leq2\operatorname{LP}_{\leq3}
 \leq2\operatorname{OPT}(G,s,c),
\]
where the last inequality holds because every globally feasible incidence
vector is feasible for the restricted LP.  Alternatively, an exact
optimization backend may return an integral minimizer \(A^{(3)}\) with
\(c(A^{(3)})=\operatorname{OPT}_{\leq3}\); no polynomial-time claim is
made for this exact integral alternative.

For either returned \(A\), put \(F_0=V\setminus A\), which contains \(s\),
and apply Lemma
\(\mathrm{lem{:}fixed\mbox{-}point\mbox{-}hedge\mbox{-}oracle}\).  If the
fixed point is \(\{s\}\), then \(F_0\) contains no hedge.  Since the graph
is finite, every hedge contains an inclusion-minimal hedge; hence \(A\)
hits all of \(\mathcal C(G,s)\).  Definition
\(\mathrm{def{:}mcid\mbox{-}objective}\) then certifies global feasibility
and identification of \(Q_s\).  For \(A^{(3)}\), restricted feasibility
gives \(\operatorname{OPT}_{\leq3}\leq\operatorname{OPT}(G,s,c)\), while
successful global feasibility gives
\[
 \operatorname{OPT}(G,s,c)
 \leq c(A^{(3)})=\operatorname{OPT}_{\leq3};
\]
thus the exact restricted optimum is globally exact.  For
\(A^{\mathrm{rnd}}\), the already established inequality gives the global
factor-two certificate.

If verification fails, the oracle returns a hedge \(F\subseteq F_0\).
For each \(v\in F\setminus\{s\}\), run the oracle on
\(F\setminus\{v\}\).  Whenever a hedge survives, replace \(F\) by the
strictly smaller fixed-point hedge returned and restart.  At most
\(|V|-1\) replacements occur, with at most \(|V|-1\) oracle calls between
replacements.  At termination no single deletion contains a hedge.  If a
proper hedge \(J\subsetneq F\) existed, choose
\(v\in F\setminus J\); both hedges contain \(s\), so \(v\neq s\), and
\(J\subseteq F\setminus\{v\}\), a contradiction.  Therefore
\(R=F\setminus\{s\}\) belongs to \(\mathcal C(G,s)\).  Since
\(F\subseteq V\setminus A\), it is missed by \(A\); because both candidate
types hit every member of \(\mathcal C_{\leq3}\), necessarily \(|R|>3\).
The final graph searches verify the hedge predicates and the failed
single-deletion searches verify minimality.

Canonical bounded enumeration, the rational LP and rounding branch,
fixed-point verification, and deletion minimization perform polynomially
many elementary encoded operations.  Clauses (i), (iii), and (iv) of
Assumption \(\mathrm{ass{:}uniform\mbox{-}poly\mbox{-}time\mbox{-}model}\)
make their composition polynomial on the canonical \(\mathrm{Fin}\,n\)
carrier; clause (ii) transfers the invariant certificate and guarantees to
the generic supplied-ordered carrier.  This proves the asserted interface
without recognizing the rank promise.

%%% FIELD proof-lem-fixed-point-hedge-oracle
For \(W\subseteq F\) with \(s\in W\), define
\[
 \mathsf A(W)=\{v\in W:v\text{ reaches }s\text{ in }D[W]\}
\]
and let \(\mathsf K(W)\) be the vertex set of the connected component of
\(s\) in \(B[W]\).  Starting from \(W_0=F\), alternate
\(W_{2j+1}=\mathsf A(W_{2j})\) and
\(W_{2j+2}=\mathsf K(W_{2j+1})\).  Both maps are deflationary and preserve
\(s\).  Each nonstationary step deletes a vertex.  Therefore, after at most
\(2|V|+2\) searches, a full two-step pass deletes nothing and the sequence
reaches a common fixed point \(F^{\star}\) satisfying
\(\mathsf A(F^{\star})=\mathsf K(F^{\star})=F^{\star}\).

If \(J\subseteq W\) is a hedge, every directed path to \(s\) inside
\(D[J]\) is also present in \(D[W]\), so \(J\subseteq\mathsf A(W)\).
Likewise every bidirected path inside the connected graph \(B[J]\) lies in
\(B[W]\), so \(J\subseteq\mathsf K(W)\).  Induction over the alternating
restrictions shows that every hedge contained in \(F\) is contained in
\(F^{\star}\).

At the common fixed point, every vertex reaches \(s\) in
\(D[F^{\star}]\), and \(B[F^{\star}]\) is connected.  If
\(F^{\star}\neq\{s\}\), Definition
\(\mathrm{def{:}singleton\mbox{-}hedges}\) therefore gives
\(F^{\star}\in\mathsf{Hdg}(G,s)\).  If
\(F^{\star}=\{s\}\), the survival implication proves that the original
\(F\) contained no hedge.  This establishes the equivalence.

Each restriction is one directed-reachability or undirected-connectivity
search on a canonically encoded induced graph.  Clause (iii) of Assumption
\(\mathrm{ass{:}uniform\mbox{-}poly\mbox{-}time\mbox{-}model}\), together
with the linear bound on the number of searches, proves polynomial runtime
on the canonical \(\mathrm{Fin}\,n\) carrier.  Clause (ii) transfers the
fixed set and the hedge predicate, both relabeling invariant, to every
finite supplied-ordered carrier.

%%% FIELD proof-lem-ugc-vertex-cover-hardness
Assume \(\mathsf{UGC}\) and fix
\(\epsilon\in\mathbb R_{>0}\).  Lemma
\(\mathrm{lem{:}ugc\mbox{-}vertex\mbox{-}cover\mbox{-}hardness\mbox{-}source}\)
gives the published \((2-\epsilon)\)-inapproximability theorem for Minimum
Vertex Cover.  Suppose, contrary to the claimed canonical specialization,
that an algorithm with that ratio ran in polynomial time on the canonical
\(\mathrm{Fin}\,n\)-encoded carrier.  Given any finite supplied-ordered
graph, clause (i) of Assumption
\(\mathrm{ass{:}uniform\mbox{-}poly\mbox{-}time\mbox{-}model}\) computes
its canonical adjacency-bit encoding, and clause (ii) transports the
returned vertex subset back along the order isomorphism with polynomial
overhead.  Vertex-cover feasibility, cardinality, and the optimum are
invariant under this relabeling.  The transported output would therefore be
a polynomial-time \((2-\epsilon)\)-approximation on the source theorem's
finite graph instances, a contradiction.  Hence the published lower bound
holds on the declared canonical carrier.  Clause (ii) is also exactly the
promised transfer back to any finite supplied-ordered carrier in later
consumers.

%%% FIELD proof-lem-rational-lp-polynomial-time
Lemma
\(\mathrm{lem{:}rational\mbox{-}lp\mbox{-}polynomial\mbox{-}time\mbox{-}source}\)
supplies an algorithm polynomial in the total binary length of a rational
linear program and guarantees an optimal basic solution whenever an optimum
exists.  For a canonically indexed program of dimensions \(d\) and total
binary length \(L\), its ordinary rational input length is bounded by a
polynomial in \(d\) and \(L\), and conversely these quantities are bounded
by its explicit encoding length.  Clause (iii) of Assumption
\(\mathrm{ass{:}uniform\mbox{-}poly\mbox{-}time\mbox{-}model}\) charges
the comparisons and rational field operations at their binary bit
complexity.  Substituting this canonical encoding into the cited algorithm
therefore gives runtime polynomial in the dimensions and \(L\), with the
same optimal-basic-solution guarantee.  Clause (ii) transports an indexed
solution along a supplied order without changing feasibility or objective
value.
